\begin{textsample}{POS Dim 3 – persona\_gpt – Score 32.00 – t167\_gpt.txt}  \label{ex:f3_pos_028}
For more than a century, Coca-Cola and Pepsi have battled for the \textbf{top} spot in the global soft drink \textbf{market}.
That rivalry has shaped everything from advertising to \textbf{product} placement, but now it \textbf{needs} to shape something far more important : their impact on the planet.
Both \textbf{companies} are among the world ’s \textbf{top} plastic polluters, flooding communities and ecosystems with single-use plastic \textbf{bottles} and \textbf{packaging}.
This plastic pollution crisis is tightly linked to the climate crisis.
Plastic is made from fossil fuels, and every step of its life cycle—from extraction and \textbf{production} to disposal—drives \textbf{greenhouse} \textbf{gas} \textbf{emissions} and environmental destruction.
As long as \textbf{companies} rely on single-use plastic, they are helping to fuel climate change and the degradation of oceans, rivers and communities.
Coca-Cola has recently committed to making at least 25 % of its \textbf{packaging} reusable or refillable by 2030.
This is far from enough to solve the \textbf{problem}, but it does set a clear benchmark in an \textbf{industry} that has been slow to move away from throwaway \textbf{packaging}.
Pepsi, by contrast, is currently at 0 % reusable or refillable \textbf{packaging} and has only pledged to set a \textbf{reuse} goal by the \textbf{end} of 2022.
In the context of the \textbf{scale} and urgency of the plastic and climate crises, that is nowhere near sufficient. \textbf{Reuse} systems—where \textbf{bottles} and containers are \textbf{designed} to be returned, washed and refilled many times—are essential to cutting plastic \textbf{waste} at the source.
They reduce the \textbf{volume} of plastic entering oceans and \textbf{waterways} and lessen the demand for new plastic \textbf{production}.
But \textbf{reuse} cannot be \textbf{scaled} up overnight.
It requires \textbf{investment} in \textbf{infrastructure}, such as collection and washing \textbf{systems}, and collaboration across \textbf{supply} \textbf{chains}, \textbf{retailers} and governments.
Global \textbf{brands} like Pepsi and Coca-Cola are uniquely positioned to drive this transformation because of their size, reach and resources.
If Coca-Cola can put a concrete \textbf{reuse} commitment on the table, Pepsi can and must do better.
Matching 25 % by 2030 would simply keep the \textbf{company} in Coca-Cola ’s shadow on an issue where leadership is urgently \textbf{needed}.
Pepsi ’s long-standing rivalry with Coca-Cola should now be about who can move fastest and furthest to phase out single-use plastics and support real \textbf{solutions}.
The upcoming May shareholder meeting is a critical moment.
Investors, \textbf{customers} and communities have an opportunity to demand that Pepsi ’s CEO commit to a meaningful, time-bound target : at least 50 % reusable or refillable \textbf{packaging} by 2030.
That kind of commitment would signal a real \textbf{shift} away from the single-use \textbf{model} that has helped make Pepsi one of the world ’s \textbf{top} plastic polluters.
Pressure from the \textbf{public} and shareholders has already begun to push major \textbf{brands} to act.
Now, it \textbf{needs} to intensify.
Before the May meeting, people can contact Pepsi ’s leadership, speak up as shareholders or \textbf{consumers}, and call for a clear, ambitious \textbf{reuse} target.
The choice facing Pepsi is simple : cling to a wasteful, polluting status quo, or finally outdo its rival where it matters most—by cutting plastic at the source and helping to tackle the climate and plastic crises together.

% matched lemmas: bottle, brand, chain, company, consumer, customer, design, emission, end, gas, greenhouse, industry, infrastructure, investment, market, model, need, packaging, problem, product, production, public, retailer, reuse, scale, shift, solution, supply, system, top, volume, waste, waterway
\end{textsample}
